% Introduction (#81, part "Abstract and introduction"; owner: 11yashiagrawal).
% Sources for numbers are listed in abstract.tex. Phase 1 numbers:
% results/metrics/fd001_test.json and fd001_test_capped125.json.

\section{Introduction}
\label{sec:introduction}

An engine that fails in service grounds the aircraft, while one that is
overhauled too early wastes useful life. Predictive maintenance estimates an
engine's remaining useful life (RUL) from its own sensor history. Most
data-driven work on this problem uses NASA's C-MAPSS
simulations~\cite{Saxena2008damage}: every training engine runs until it fails,
and every test engine stops at an unknown point before failure.

Phase~1 of this project built feature-based baselines on FD001, the subset
with one operating condition and one fault mode, under a protocol fixed for
both phases. The protocol splits by engine, refits every preprocessing step
inside each training fold, and scores each official test engine once, at its
last cycle. The best Phase~1 model, a Random Forest on 155 engineered rolling
features, reached a test RMSE of 27.55 cycles on the uncapped target. Capping
the training target at 125 cycles lowered it to 16.21. Two questions were left
open. First, the cap had been taken from the literature and was already known
to work on the test set, so it was not a fair choice. Second, fixed rolling
summaries may miss degradation patterns that a model reading the whole
sequence could find~\cite{Zheng2017long}.

Phase~2 addresses both questions. Published deep-learning results on FD001
range from about 11 to 17 cycles, but they use different caps, test labels,
windows and numbers of runs (Section~\ref{sec:lit_dl}), so they cannot tell us
whether a sequence model beats a strong feature baseline. We therefore compare
the two families under one protocol and one target, and ask three questions:
\begin{enumerate}
\item Do deep sequence models beat machine learning on engineered features,
  when both train on the same capped target and are selected with the same
  metric?
\item Which way of reading the window works best on FD001: recurrence (LSTM),
  dilated convolution (temporal CNN) or attention (Transformer)?
\item Does fusing a sequence encoder with the engineered features add value
  beyond the better of the two, under a test fixed in advance?
\end{enumerate}

The contributions of this phase are:
\begin{itemize}
\item a rule that chooses the RUL cap from training engines alone (health
  index, hinge-onset RUL, 90th percentile), together with the selection metric
  (RMSE on cycles with true RUL $\leq K$) used for all tuning;
\item a shared encoder--head contract under which the LSTM, temporal CNN and
  Transformer differ only in their encoder, and machine-learning baselines
  retuned on the same target and metric;
\item two hybrids, gated fusion and a two-stage model, judged by a value test
  set in advance: a paired bootstrap, a Holm-corrected Wilcoxon test, or lower
  variance across seeds;
\item ablations, technical validation and interpretability analyses, and a
  single run on the official test set after the model freeze.
\end{itemize}

Section~\ref{sec:lit_dl} reviews deep learning on C-MAPSS.
Section~\ref{sec:target_selection} chooses the cap and the selection metric,
Section~\ref{sec:dl_dataset} builds the sequence dataset,
Section~\ref{sec:theory} explains the models and Section~\ref{sec:architecture}
the shared architecture. The hybrid models, ablations, validation,
interpretability and the official results follow, and the last section
concludes.
